\documentclass[10pt,a4paper]{amsart}
\usepackage{amsmath,amssymb,amsthm,graphicx}
% For commutative diagrams we use Xy-pic
\usepackage[all]{xy}

% We use 2cell for 2-commutative diagrams.
\xyoption{2cell}
% Theorem environments.
%
\theoremstyle{plain}
\newtheorem{theorem}[subsection]{Theorem}
\newtheorem{proposition}[subsection]{Proposition}
\newtheorem{lemma}[subsection]{Lemma}

\theoremstyle{definition}
\newtheorem{definition}[subsection]{Definition}
\newtheorem{example}[subsection]{Example}
\newtheorem{exercise}[subsection]{Exercise}
\newtheorem{situation}[subsection]{Situation}

\theoremstyle{remark}
\newtheorem{remark}[subsection]{Remark}
\newtheorem{remarks}[subsection]{Remarks}

\numberwithin{equation}{subsection}

% Macros
%
\def\lim{\mathop{\mathrm{lim}}\nolimits}
\def\colim{\mathop{\mathrm{colim}}\nolimits}
\def\Spec{\mathop{\mathrm{Spec}}}
\def\Hom{\mathop{\mathrm{Hom}}\nolimits}
\def\Ext{\mathop{\mathrm{Ext}}\nolimits}
\def\SheafHom{\mathop{\mathcal{H}\!\mathit{om}}\nolimits}
\def\SheafExt{\mathop{\mathcal{E}\!\mathit{xt}}\nolimits}
\def\Sch{\mathit{Sch}}
\def\Mor{\mathop{\mathrm{Mor}}\nolimits}
\def\Ob{\mathop{\mathrm{Ob}}\nolimits}
\def\Sh{\mathop{\mathit{Sh}}\nolimits}
\def\NL{\mathop{N\!L}\nolimits}
\def\CH{\mathop{\mathrm{CH}}\nolimits}
\def\proetale{{pro\text{-}\acute{e}tale}}
\def\etale{{\acute{e}tale}}
\def\QCoh{\mathit{QCoh}}
\def\Ker{\mathop{\mathrm{Ker}}}
\def\Im{\mathop{\mathrm{Im}}}
\def\Coker{\mathop{\mathrm{Coker}}}
\def\Coim{\mathop{\mathrm{Coim}}}

% Boxtimes
%
\DeclareMathSymbol{\boxtimes}{\mathbin}{AMSa}{"02}

%
% Macros for moduli stacks/spaces
%
\def\QCohstack{\mathcal{QC}\!\mathit{oh}}
\def\Cohstack{\mathcal{C}\!\mathit{oh}}
\def\Spacesstack{\mathcal{S}\!\mathit{paces}}
\def\Quotfunctor{\mathrm{Quot}}
\def\Hilbfunctor{\mathrm{Hilb}}
\def\Curvesstack{\mathcal{C}\!\mathit{urves}}
\def\Polarizedstack{\mathcal{P}\!\mathit{olarized}}
\def\Complexesstack{\mathcal{C}\!\mathit{omplexes}}
% \Pic is the operator that assigns to X its picard group, usage \Pic(X)
% \Picardstack_{X/B} denotes the Picard stack of X over B
% \Picardfunctor_{X/B} denotes the Picard functor of X over B
\def\Pic{\mathop{\mathrm{Pic}}\nolimits}
\def\Picardstack{\mathcal{P}\!\mathit{ic}}
\def\Picardfunctor{\mathrm{Pic}}
\def\Deformationcategory{\mathcal{D}\!\mathit{ef}}
\usepackage{cleveref}
\crefname{originallemma}{original lemma}{original lemmas}
\Crefname{originallemma}{Original lemma}{Original lemmas}
\makeatletter
\newlabel{lemma-localized-chern-pre}{{0F9K}{source}{}{sourcepointer-1}{}}
\newlabel{lemma-localized-chern-pre@cref}{{[originallemma][0][]0F9K}{[1][0][]source}{}{}{}}
\makeatother
\crefname{originallemma}{original lemma}{original lemmas}
\Crefname{originallemma}{Original lemma}{Original lemmas}
\makeatletter
\newlabel{lemma-gysin-at-infty}{{0F9J}{source}{}{sourcepointer-2}{}}
\newlabel{lemma-gysin-at-infty@cref}{{[originallemma][0][]0F9J}{[1][0][]source}{}{}{}}
\makeatother
\crefname{originallemma}{original lemma}{original lemmas}
\Crefname{originallemma}{Original lemma}{Original lemmas}
\makeatletter
\newlabel{lemma-bivariant-zero}{{02UC}{source}{}{sourcepointer-3}{}}
\newlabel{lemma-bivariant-zero@cref}{{[originallemma][0][]02UC}{[1][0][]source}{}{}{}}
\makeatother
\crefname{originallemma}{original lemma}{original lemmas}
\Crefname{originallemma}{Original lemma}{Original lemmas}
\makeatletter
\newlabel{lemma-silly}{{0F9H}{source}{}{sourcepointer-4}{}}
\newlabel{lemma-silly@cref}{{[originallemma][0][]0F9H}{[1][0][]source}{}{}{}}
\makeatother
\crefname{originallemma}{original lemma}{original lemmas}
\Crefname{originallemma}{Original lemma}{Original lemmas}
\makeatletter
\newlabel{lemma-silly-silly}{{0F9I}{source}{}{sourcepointer-5}{}}
\newlabel{lemma-silly-silly@cref}{{[originallemma][0][]0F9I}{[1][0][]source}{}{}{}}
\makeatother
\crefname{originallemma}{original lemma}{original lemmas}
\Crefname{originallemma}{Original lemma}{Original lemmas}
\makeatletter
\newlabel{lemma-homomorphism-pre}{{0FAR}{source}{}{sourcepointer-6}{}}
\newlabel{lemma-homomorphism-pre@cref}{{[originallemma][0][]0FAR}{[1][0][]source}{}{}{}}
\makeatother
\title{Intertwining of the Gysin-at-infinity class with the localized p-th Chern class under the further hypothesis that the rank of the quotient is less than p}
\author{The Stacks Project Authors}
\date{stacks-project:chow.tex:local-snapshot-20261012; GNU FDL 1.2 or later}
\begin{document}
\maketitle
\noindent\textit{Editorial note.} All delivered mathematical text is the author's own, copied from the pinned Stacks Project chapter file chow.tex; no statement or proof step was rewritten, reordered or removed.\par\medskip
\noindent\textit{Editorial note.} Preamble provenance: the pinned chapter file opens with an input of the shared preamble and carries no preamble text of its own, so the project's single Stacks-wide preamble staged at sources/stacks-chow/preamble.tex is excerpted verbatim as two spans, lines 14-18 (Xy-pic setup) and lines 154-217 (the native theorem declarations and mathematical macros). Nothing was added, removed or redefined; the wrapper class and the amsmath/amssymb/amsthm/graphicx base packages are supplied by the assembler.\par\medskip
\noindent\textit{Editorial note.} The statement inherits the full setup of the same chapter result labelled lemma-localized-chern-pre (retained as a tagged navigation pointer) together with the additional rank hypothesis rank(Q|\_T) < p; the arguments of the lemma are the two chain computations written out inside the counted span.\par\medskip
\noindent\textit{Editorial note.} Canonical location of the counted proposition: Stacks tag 0FAV (sources/stacks-chow/pinned-original-source.tex line 10353).\par\medskip
\noindent\textit{Allowed standard prerequisites (editorial).}\begin{itemize}
\item The localized Chern class setup and the class c'\_p(Q). Same chapter, Lemma labelled lemma-localized-chern-pre (source line 10204), whose setup the statement inherits, invoked at statement line 10354 and proof line 10366. Fix the objects Q, Z, T, W, b and the classes c\_p(Q), c'\_p(Q) and the complement assumption rank(Q|\_T) < p.
\item The Gysin-at-infinity bivariant class C and its defining property. Same chapter, Lemma labelled lemma-gysin-at-infty (source line 10026), invoked at statement line 10357. Provides the class C with C cap alpha = i\_infinity\textasciicircum{}* beta for a suitable beta, which is the third equality of the displayed chain.
\item Two intermediate computations identifying the restricted classes. Same chapter, Lemmas labelled lemma-silly (source line 9695), lemma-silly-silly (source line 9792), lemma-homomorphism-pre (source line 10132) and lemma-bivariant-zero (source line 6585), invoked at proof lines 10368, 10395, 10400 and 10402. Supply the individual steps of the chain of eight equalities, namely the description of c'\_p(Q) through E, the comparison on i\_infinity\textasciicircum{}* beta, the commutation with i\_infinity\textasciicircum{}* and the reduction to testing against cycle classes.
\end{itemize}\medskip
\ifdefined\hypertarget\hypertarget{sourcepointer-1}{}\fi
\noindent\textit{Source reference.} Original lemma 0F9K: sources/stacks-chow/pinned-original-source.tex, source lines 10204--10204. Same-chapter pointer into the author environment lemma at source line 10204; invoked at proof line 10366 of the counted proof. It is retained only as a navigation pointer; every load-bearing step is inside the delivered span.\par\medskip
\ifdefined\hypertarget\hypertarget{sourcepointer-2}{}\fi
\noindent\textit{Source reference.} Original lemma 0F9J: sources/stacks-chow/pinned-original-source.tex, source lines 10026--10026. Same-chapter pointer into the author environment lemma at source line 10026; invoked at proof line None of the counted proof. It is retained only as a navigation pointer; every load-bearing step is inside the delivered span.\par\medskip
\ifdefined\hypertarget\hypertarget{sourcepointer-3}{}\fi
\noindent\textit{Source reference.} Original lemma 02UC: sources/stacks-chow/pinned-original-source.tex, source lines 6585--6585. Same-chapter pointer into the author environment lemma at source line 6585; invoked at proof line 10368 of the counted proof. It is retained only as a navigation pointer; every load-bearing step is inside the delivered span.\par\medskip
\ifdefined\hypertarget\hypertarget{sourcepointer-4}{}\fi
\noindent\textit{Source reference.} Original lemma 0F9H: sources/stacks-chow/pinned-original-source.tex, source lines 9695--9695. Same-chapter pointer into the author environment lemma at source line 9695; invoked at proof line 10395 of the counted proof. It is retained only as a navigation pointer; every load-bearing step is inside the delivered span.\par\medskip
\ifdefined\hypertarget\hypertarget{sourcepointer-5}{}\fi
\noindent\textit{Source reference.} Original lemma 0F9I: sources/stacks-chow/pinned-original-source.tex, source lines 9792--9792. Same-chapter pointer into the author environment lemma at source line 9792; invoked at proof line 10400 of the counted proof. It is retained only as a navigation pointer; every load-bearing step is inside the delivered span.\par\medskip
\ifdefined\hypertarget\hypertarget{sourcepointer-6}{}\fi
\noindent\textit{Source reference.} Original lemma 0FAR: sources/stacks-chow/pinned-original-source.tex, source lines 10132--10132. Same-chapter pointer into the author environment lemma at source line 10132; invoked at proof line 10402 of the counted proof. It is retained only as a navigation pointer; every load-bearing step is inside the delivered span.\par\medskip
\begin{lemma}
\label{lemma-homomorphism}
In Lemma \ref{lemma-localized-chern-pre} assume $Q|_T$ is isomorphic
to a finite locally free $\mathcal{O}_T$-module of rank $< p$.
Denote $C \in A^0(W_\infty \to X)$ the class of
Lemma \ref{lemma-gysin-at-infty}. Then
$$
C \circ c_p(Q|_{X \times \{0\}}) =
C \circ (Z \to X)_* \circ c'_p(Q) = c_p(Q|_{W_\infty}) \circ C
$$
\end{lemma}

\begin{proof}
The first equality holds because $c_p(Q|_{X \times \{0\}}) =
(Z \to X)_* \circ c'_p(Q)$ by Lemma \ref{lemma-localized-chern-pre}.
We may prove the second equality one cycle class at a time
(see Lemma \ref{lemma-bivariant-zero}). Since the construction of
the bivariant classes in the lemma is compatible with base change,
we may assume we have some $\alpha \in \CH_k(X)$ and we have to show that
$C \cap (Z \to X)_*(c'_p(Q) \cap \alpha) =
c_p(Q|_{W_\infty}) \cap C \cap \alpha$. Observe that
\begin{align*}
C \cap (Z \to X)_*(c'_p(Q) \cap \alpha)
& =
C \cap (Z \to X)_* (E \to Z)_*(c'_p(Q|_E) \cap C \cap \alpha) \\
& =
C \cap (W_\infty \to X)_*(E \to W_\infty)_*(c'_p(Q|_E) \cap C \cap \alpha) \\
& =
C \cap (W_\infty \to X)_*(E \to W_\infty)_*(c'_p(Q|_E) \cap i_\infty^*\beta) \\
& =
C \cap (W_\infty \to X)_*(c_p(Q|_{W_\infty}) \cap i_\infty^*\beta) \\
& =
C \cap (W_\infty \to X)_*i_\infty^*(c_p(Q) \cap \beta) \\
& =
i_\infty^*(c_p(Q) \cap \beta) \\
& =
c_p(Q|_{W_\infty}) \cap i_\infty^*\beta \\
& =
c_p(Q|_{W_\infty}) \cap C \cap \alpha
\end{align*}
as desired. For the first equality we used that
$c'_p(Q) = (E \to Z)_* \circ c'_p(Q|_E) \circ C$ where $E \subset W_\infty$
is the inverse image of $Z$ and $c'_p(Q|_E)$ is the class constructed
in Lemma \ref{lemma-silly}. The second equality is just the statement
that $E \to Z \to X$ is equal to $E \to W_\infty \to X$.
For the third equality we choose $\beta \in \CH_{k + 1}(W)$ whose restriction to
$b^{-1}(\mathbf{A}^1_X)$ is the flat pullback of $\alpha$ so that
$C \cap \alpha = i_\infty^*\beta$ by construction. The fourth equality is
Lemma \ref{lemma-silly-silly}. The fifth equality is the fact that
$c_p(Q)$ is a bivariant class and hence commutes with $i_\infty^*$.
The sixth equality is Lemma \ref{lemma-homomorphism-pre}.
The seventh uses again that $c_p(Q)$ is a bivariant class.
The final holds as $C \cap \alpha = i_\infty^*\beta$.
\end{proof}

\end{document}
